\documentclass[11pt]{article}

\usepackage[a4paper,margin=1.4cm]{geometry}
\usepackage{multicol}
\usepackage{xcolor}
\usepackage{graphicx}
\usepackage{tcolorbox}
\usepackage{listings}
\usepackage{titlesec}
\usepackage{parskip}
\usepackage{amsmath}

\setlength{\columnsep}{1cm}

\definecolor{mainblue}{RGB}{30,144,255}
\definecolor{lightblue}{RGB}{220,240,255}
\definecolor{lightgray}{RGB}{245,245,245}
\definecolor{dark}{RGB}{40,40,40}

\lstset{
    basicstyle=\ttfamily\small,
    backgroundcolor=\color{lightgray},
    frame=single,
    breaklines=true
}

% Blocks
\newtcolorbox{missionbox}{
    colback=mainblue!10,
    colframe=mainblue,
    title=\textbf{Mission}
}

\newtcolorbox{theorybox}{
    colback=white,
    colframe=black!60,
    title=\textbf{Concept}
}

\newtcolorbox{examplebox}{
    colback=lightgray,
    colframe=black,
    title=\textbf{Observe}
}

\newtcolorbox{trybox}{
    colback=green!5,
    colframe=green!60!black,
    title=\textbf{Try it yourself}
}

\newtcolorbox{bossbox}{
    colback=purple!10,
    colframe=purple,
    title=\textbf{Boss Level}
}

\begin{document}

\begin{center}
{\Huge \textbf{CODE QUEST}}\\
Algorithms World\\
Level 1 • Issue 1\\[10pt]

{\huge \textbf{CALCULATORUL NU ȘTIE SĂ NUMERE}}\\
O lecție de la variabile până la pătrate
\end{center}

\vspace{0.5cm}

\begin{multicols}{2}

% Mission
\begin{missionbox}
Calculatorul NU știe să numere.

El știe doar:
\begin{itemize}
\item să țină minte numere
\item să facă operații simple (+, -)
\item să sară între instrucțiuni
\end{itemize}

Dacă îl învățăm să numere, putem construi lumi — exact ca în Minecraft.
\end{missionbox}

% Variables
\section*{Memoria}

\begin{theorybox}
Variabilele sunt zone de memorie în care calculatorul păstrează valori.

Tipul cel mai folosit:
\textbf{int} = număr întreg
\end{theorybox}

\begin{examplebox}
\begin{lstlisting}
int a = 1;
int b = 10;
\end{lstlisting}
\end{examplebox}

\begin{trybox}
Declară trei variabile:
x = 5, y = 12, scor = 0
\end{trybox}

% Operators
\section*{Aritmetica}

\begin{examplebox}
\begin{lstlisting}
i = i + 1;
i += 1;
i++;
\end{lstlisting}
Toate înseamnă: crește i cu 1
\end{examplebox}

\begin{trybox}
Dacă i = 7, calculează:
i++; i+=3; i--;
\end{trybox}

% Control flow
\section*{Salturi}

\begin{theorybox}
Calculatorul poate să sară între instrucțiuni.
\end{theorybox}

\begin{examplebox}
\begin{lstlisting}
cmp i, 10
jle inside
jmp out
\end{lstlisting}
\end{examplebox}

% While
\section*{Repetiția}

\begin{examplebox}
\begin{lstlisting}
int i = 1;
while(i <= 10)
{
    cout << i << " ";
    i++;
}
\end{lstlisting}
\end{examplebox}

\begin{trybox}
Modifică programul pentru 1 → 20
\end{trybox}

% Key idea
\section*{Cum numărăm?}

\begin{theorybox}
3 lucruri:
\begin{itemize}
\item Start: i = 1
\item Stop: i <= 10
\item Pas: i++
\end{itemize}
\end{theorybox}

% For
\section*{Structura for}

\begin{examplebox}
\begin{lstlisting}
for(int i = 1; i <= 10; i++)
{
    cout << i;
}
\end{lstlisting}
\end{examplebox}

\begin{trybox}
Afișează doar numerele pare
\end{trybox}

% Fizz
\section*{Condiții}

\begin{examplebox}
\begin{lstlisting}
for(int i = 1; i <= 10; i++)
{
    if(i % 2 == 0)
        cout << i;
}
\end{lstlisting}
\end{examplebox}

\begin{trybox}
Afișează numerele divizibile cu 3
\end{trybox}

% Double for
\section*{Pătrate}

\begin{examplebox}
\begin{lstlisting}
for(int i = 1; i <= 5; i++)
{
    for(int j = 1; j <= 5; j++)
    {
        cout << "#";
    }
    cout << endl;
}
\end{lstlisting}
\end{examplebox}

\begin{trybox}
Desenează un pătrat de dimensiune n
\end{trybox}

% Diagonals
\section*{Diagonale}

\begin{examplebox}
\begin{lstlisting}
if(i == j)
    cout << "*";
else
    cout << ".";
\end{lstlisting}
\end{examplebox}

\begin{trybox}
Desenează diagonala secundară
\end{trybox}

% Triangle
\section*{Triunghiuri}

\begin{examplebox}
\begin{lstlisting}
for(int i = 1; i <= n; i++)
{
    for(int j = 1; j <= i; j++)
        cout << "#";
    cout << endl;
}
\end{lstlisting}
\end{examplebox}

% Boss
\begin{bossbox}
Se citește n.

Desenează un pătrat n x n:
\begin{itemize}
\item \# margine
\item * diagonala principală
\item + diagonala secundară
\item . interior
\end{itemize}
\end{bossbox}

\end{multicols}

\end{document}
